\reading{IM-1}{Factor Theory}{STRIKE FIRST}{5}

\srcnote{Andrew Ang, \textit{Asset Management: A Systematic Approach to Factor
Investing} (New York, NY: Oxford University Press, 2014), Chapter 6.
Printed as Chapter 1 in the GARP source volume.}

\begin{imkeybox}[title={The six objectives in this reading}]
\begin{enumerate}[label=\textbf{\alph*.}]
\item Describe factors that impact asset prices and explain the theory of factor risk premiums.
\item Discuss the capital asset pricing model (CAPM) including its assumptions and explain how factor risk is addressed in the CAPM.
\item Explain the implications of using the CAPM to value assets, including equilibrium and optimal holdings, exposure to factor risk, its treatment of diversification benefits, and shortcomings of the CAPM.
\item Describe multifactor models and compare and contrast multifactor models to the CAPM.
\item Explain how stochastic discount factors are created and apply them in the valuation of assets.
\item Describe efficient market theory and explain how markets can be inefficient.
\end{enumerate}
\end{imkeybox}

% ------------------------------------------------------------------
\lo{IM-1 a}{Describe factors that impact asset prices and explain the theory of factor risk premiums.}

\subsection{Factors that impact asset prices}

There are three similarities between food and assets:

\begin{itemize}
\item Factors matter, not assets.
\item Assets are bundles of factors.
\item Different investors need different risk factors.
\end{itemize}

There is one difference, however, between factors and nutrients. Nutrients are
inherently good for you. \term{Factor risks are bad.}

\begin{center}
\begin{tikzpicture}[
  node distance=5mm,
  bx/.style={draw=FRMNavy,fill=PaleNavy,rounded corners=1pt,inner sep=4pt,
             font=\scriptsize\sffamily,text=FRMNavy,minimum width=17mm,align=center},
  hd/.style={draw=FRMNavy,fill=FRMNavy,rounded corners=1pt,inner sep=4pt,
             font=\scriptsize\sffamily\bfseries,text=white,minimum width=22mm,align=center},
  ar/.style={-{Stealth[length=4pt]},FRMGold,line width=1pt}]

% left group
\node[hd] (food) at (0,0) {Food};
\node[hd,below=9mm of food] (nutr) {Nutrients};
\draw[ar] (food) -- (nutr);
\node[bx,below left=9mm and 10mm of nutr.south] (n1) {Fibre};
\node[bx,below=9mm of nutr] (n2) {Fat};
\node[bx,below right=9mm and 10mm of nutr.south] (n3) {Protein};
\draw[ar] (nutr) -- (n1); \draw[ar] (nutr) -- (n2); \draw[ar] (nutr) -- (n3);

% right group
\node[hd] (asset) at (7.4,0) {Asset};
\node[hd,below=9mm of asset] (rf) {Risk factors};
\draw[ar] (asset) -- (rf);
\node[bx,below left=9mm and 10mm of rf.south] (r1) {Market\\risk};
\node[bx,below=9mm of rf] (r2) {Credit\\risk};
\node[bx,below right=9mm and 10mm of rf.south] (r3) {Liquidity\\risk};
\draw[ar] (rf) -- (r1); \draw[ar] (rf) -- (r2); \draw[ar] (rf) -- (r3);

% mapping
\draw[{Stealth[length=5pt]}-{Stealth[length=5pt]},FRMGold,line width=1pt]
  (nutr.east) -- node[above,font=\scriptsize\sffamily,fill=white,inner sep=1.5pt,text=FRMGold]{same idea}
  node[below,font=\scriptsize\sffamily,fill=white,inner sep=1.5pt,text=FRMRed]{opposite sign} (rf.west);
\end{tikzpicture}
\end{center}
\figcap{The analogy and the place where it breaks. You do not eat nutrients
directly, and you do not hold factors directly; both are consumed through the
bundles that contain them. The mapping is exact until you reach the sign:
nutrients are good and you pay for them, whereas factor risks are bad and you
are paid to bear them.}

\subsection{The theory of factor risk premiums}

The factor theory is trying to explain the phenomenon:

\begin{imkeybox}
Investors exposed to losses during bad times are compensated by risk premiums in
good times.
\end{imkeybox}

% ------------------------------------------------------------------
\lo{IM-1 b}{Discuss the capital asset pricing model (CAPM) including its assumptions and explain how factor risk is addressed in the CAPM.}

The basic intuition of the CAPM still holds true:

\begin{itemize}
\item That the factors underlying the assets determine asset risk premiums and
that these risk premiums are compensation for investors' losses during bad times.
\item It predicts that asset risk premiums depend only on the asset's beta and
there is only one factor that matters, the market portfolio.
\end{itemize}

\begin{imfmlbox}
\[ E(R_P) = R_f + \beta_P\left[E(R_M) - R_f\right] \]
\vk{$E(R_P)$ is the expected return on the portfolio; $R_f$ the risk-free rate;
$\beta_P$ the portfolio's sensitivity to the single priced factor; and
$E(R_M)-R_f$ the market risk premium.}
\end{imfmlbox}

\subsection{The seven assumptions behind the CAPM}

\begin{enumerate}
\item Investors only have financial wealth.
\item Investors have mean-variance utility.
\item Investors have single period investment horizon.
\item Investors have homogeneous (identical) expectations.
\item No taxes or transactions costs.
\item All investors are price takers.
\item Information is free and available to everyone.
\end{enumerate}

\begin{imtrapbox}[title={Assumption 5 is about taxes, not just friction}]
The source notes give assumption 5 as ``markets are frictionless.'' GARP's
wording is \term{no taxes or transactions costs}, and the difference is
examinable: Ang treats taxes as a systematic factor in their own right, which the
word ``frictionless'' does not convey. When this assumption is relaxed, two
distinct premiums appear, a tax effect and a liquidity effect.
\end{imtrapbox}

% ------------------------------------------------------------------
\lo{IM-1 c}{Explain the implications of using the CAPM to value assets, including equilibrium and optimal holdings, exposure to factor risk, its treatment of diversification benefits, and shortcomings of the CAPM.}

\subsection{Lessons learnt from the CAPM}

\begin{enumerate}
\item \term{Don't hold an individual asset, hold the factor} (the market risk premium).
\item Each investor has his \term{own optimal exposure} of factor risk.
\item The average investor holds the market (MVE $=$ market portfolio).
\item The factor risk premium has an economic story.
\item \term{Risk is factor exposure} (beta matters how much the risk is to be taken).
\item Assets paying off in bad times have low risk premiums; in some
circumstances they are also defined as valuable assets such as gold.
\end{enumerate}

\begin{imfmlbox}
\[ E(R_m) - R_f \;=\; \bar{\gamma}\,\sigma_m^{2} \]
\vk{$\bar{\gamma}$ is the risk aversion of the ``average'' investor --- the
wealth-weighted average degree of risk aversion across all investors --- and
$\sigma_m^{2}$ is the variance of the market portfolio.}
\end{imfmlbox}

\noindent which indicates that the risk premium should be compensated by average
risk aversion or the variance of the market.

\begin{imtrapbox}[breakable=false,title={A symbol the source notes lost}]
The notes render this equation as $E(R_m)-R_f = {}^{-}\sigma^2$, with the
$\bar{\gamma}$ reduced to a stray macron. Restored above from the GARP source
(Equation 1.1). Both drivers matter and each is a separate lever: the premium
rises when the average investor becomes \emph{more} risk averse, and it rises
when the market becomes \emph{more} volatile. A distractor that raises one and
lowers the other is the natural corruption here.
\end{imtrapbox}

\subsection{Shortcomings of the CAPM}

The seven assumptions above are the shortcomings: each one, when relaxed,
generates an additional risk premium that the CAPM cannot price. Deviations from
the CAPM are expected to be largest in \term{illiquid, inefficient markets},
where the perfect-information assumption fails hardest.

% ------------------------------------------------------------------
\lo{IM-1 d}{Describe multifactor models and compare and contrast multifactor models to the CAPM.}

\begin{itemize}
\item Multifactor models recognize that \term{bad times can be defined more
broadly} than just bad returns on the market portfolio.
\item The first multifactor model was the \term{arbitrage pricing theory (APT)},
developed by Stephen Ross (1976).
\item In equilibrium, investors must be compensated for bearing these multiple
sources of factor risk.
\item While the CAPM captures the notion of bad times solely by means of low
returns of the market portfolio, \term{each factor in a multifactor model
provides its own definition of bad times}.
\end{itemize}

\gapbadge
\srcnote{The notes state the difference between the two models but do not carry
the comparison itself. The table below is written from the GARP curriculum
(Ang, Chapter 1, \textit{Multifactor Model Lessons}), which sets the same six
lessons side by side --- which is precisely what this objective's
``compare and contrast'' asks for.}

\begin{imgapbox}[title={What the objective asks for}]
Not ``what is a multifactor model'' but \emph{what changes and what does not}
when you move from one factor to many. Five of the six lessons survive
essentially unchanged; only the wording of Lessons 1, 4 and 5 shifts.
\end{imgapbox}

\renewcommand{\arraystretch}{1.25}
\noindent\begin{tabularx}{\textwidth}{@{}p{1.5cm}XX@{}}
\toprule
& \textbf{\sffamily\small CAPM (market factor)} & \textbf{\sffamily\small Multifactor models} \\
\midrule
\textbf{Lesson 1} & Diversification works. The market diversifies away idiosyncratic risk. & Diversification works. The \emph{tradeable version of a factor} diversifies away idiosyncratic risk. \\
\textbf{Lesson 2} & Each investor has her own optimal exposure of the market portfolio. & Each investor has her own optimal exposure of \emph{each} factor risk. \\
\textbf{Lesson 3} & The average investor holds the market. & The average investor holds the market. \\
\textbf{Lesson 4} & The market factor is priced in equilibrium under the CAPM assumptions. & Risk premiums exist for each factor assuming \emph{no arbitrage or equilibrium}. \\
\textbf{Lesson 5} & Risk of an asset is measured by the CAPM beta. & Risk of an asset is measured in terms of the \emph{factor exposures (factor betas)} of that asset. \\
\textbf{Lesson 6} & Assets paying off in bad times \emph{when the market return is low} are attractive, and these assets have low risk premiums. & Assets paying off in bad times are attractive, and these assets have low risk premiums. \\
\bottomrule
\end{tabularx}
\renewcommand{\arraystretch}{1}
\figcap{The six CAPM lessons and their multifactor counterparts. Lesson 3 is
identical in both columns, and it is the one candidates most often assume must
have changed. Lesson 4 is the sharpest distinction: the CAPM needs full
equilibrium, whereas a multifactor risk premium needs only no arbitrage.}

% ------------------------------------------------------------------
\lo{IM-1 e}{Explain how stochastic discount factors are created and apply them in the valuation of assets.}

\subsection{Pricing kernels}

To capture the composite bad times over multiple factors, the new asset pricing
approach uses the notion of a \term{pricing kernel}. This is also called a
\term{stochastic discount factor (SDF)}.

Just as the CAPM gives rise to assets having betas with respect to the market,
multiple factors in the SDF give rise to a multi-beta relation for an asset's
risk premium:

\begin{imfmlbox}
\[ E(r_i) = r_f + \beta_{i,1}E(f_1) + \beta_{i,2}E(f_2) + \cdots + \beta_{i,K}E(f_K) \]
\vk{$\beta_{i,k}$ is the beta of asset $i$ with respect to factor $k$, and
$E(f_k)$ is the risk premium of factor $k$.}
\end{imfmlbox}

For macro factors, $f_1$ could be inflation and $f_2$ could be economic growth,
for example. Bad times are characterized by times of high inflation, low economic
growth, or both.

\gapbadge
\srcnote{The objective has two verbs --- \emph{created} and \emph{applied in
valuation} --- and the notes serve neither. They give the end result, the
multi-beta relation above, without the SDF that produces it or the pricing
equation that uses it. Both halves are written below from the GARP curriculum
(Ang, Chapter 1, Equations 1.3 to 1.8).}

\begin{imgapbox}[title={What the objective asks for}]
How $m$ is \emph{built} (it is a function of the factors, and the CAPM is the
special case where that function is linear in the market return), and how $m$ is
\emph{used} to price an asset (multiply the payoff by $m$ and take an
expectation, instead of dividing the payoff by one plus a discount rate).
\end{imgapbox}

\subsection{How the SDF is created}

We denote the SDF as $m$. It is an index of bad times, indexed by many different
factors and states of nature. The CAPM is the special case in which $m$ is
\emph{linear in the market return}:

\begin{imfmlbox}
\[ m = a + b \times r_m \]
\vk{$a$ and $b$ are constants. The coefficient $b$ is \term{negative}: low values
of the SDF correspond to bad times, which in the CAPM are given by low returns of
the market.}
\end{imfmlbox}

Generalising, the SDF is allowed to depend on a vector of factors
$F = (f_1, f_2, \ldots, f_K)$, so that $m = m(F)$, where each of the $K$ factors
itself defines a different kind of bad time. A second advantage follows: the
CAPM restricts $m$ to be linear, whereas the world is nonlinear, and the general
form can carry that nonlinearity.

\begin{imtrapbox}[title={Sign and requirement, both corruptible}]
$b$ is negative, not positive. And the requirement for an SDF to exist is only
\term{no arbitrage}; it is the CAPM, and other specific forms of $m$, that
require the additional onerous assumptions. A distractor that reverses this ---
claiming the SDF needs equilibrium while the CAPM needs only no arbitrage ---
is the standard corruption.
\end{imtrapbox}

\subsection{Applying the SDF in valuation}

\begin{center}
\begin{tikzpicture}[
  bx/.style={draw=FRMRule,fill=white,rounded corners=1pt,inner sep=5pt,
             font=\scriptsize\sffamily,align=center,text width=38mm},
  hd/.style={draw=FRMNavy,fill=FRMNavy,rounded corners=1pt,inner sep=4pt,
             font=\scriptsize\sffamily\bfseries,text=white,align=center,text width=38mm},
  ar/.style={-{Stealth[length=4pt]},FRMGold,line width=1pt}]
\node[hd] (a) at (0,0) {Discount-rate model};
\node[bx,below=5mm of a] (a2) {$P_i = E\!\left[\dfrac{\text{payoff}_i}{1+E(r_i)}\right]$\\[3pt]
  discount rate from the CAPM:\\ $E(r_i)=r_f+\beta_i[E(r_m)-r_f]$};
\draw[ar] (a) -- (a2);
\node[hd] (b) at (7.0,0) {SDF model};
\node[bx,below=5mm of b] (b2) {$P_i = E\!\left[m \times \text{payoff}_i\right]$\\[3pt]
  the discounting is done \emph{inside}\\ the expectation, state by state};
\draw[ar] (b) -- (b2);
\draw[{Stealth[length=5pt]}-{Stealth[length=5pt]},FRMGold,line width=1pt]
  (a2.east) -- node[above,font=\scriptsize\sffamily,fill=white,inner sep=1.5pt,text=FRMGold]
  {same price} (b2.west);
\end{tikzpicture}
\end{center}
\figcap{Two routes to one price. The traditional model divides by a single
discount rate that is constant across states; the SDF model multiplies by $m$,
which varies across states, so bad states are weighted more heavily. Where the
CAPM works, both give the same answer --- the SDF route simply does not need the
CAPM's assumptions to get there.}

Dividing both sides of the pricing equation by the asset's current price gives a
relation that must hold for every asset:

\begin{imfmlbox}
\[ 1 = E\!\left[m \times (1+r_i)\right] \]
A special case occurs when the payoffs are constant, which gives a risk-free
asset, so the price of a risk-free bond is simply
$\dfrac{1}{1+r_f} = E\!\left[m \times 1\right] = E(m)$.
\end{imfmlbox}

The risk premium can then be written in a relation closely parallel to the
security market line:

\begin{imfmlbox}
\[ E(r_i) - r_f \;=\; \underbrace{\frac{\operatorname{cov}(r_i,m)}{\operatorname{var}(m)}}_{\textstyle \beta_{i,m}}
   \;\times\; \underbrace{\left(-\frac{\operatorname{var}(m)}{E(m)}\right)}_{\textstyle \lambda_m} \]
\vk{$\beta_{i,m}$ is the beta of the asset with respect to the SDF, and
$\lambda_m$ is the price of ``bad times'' risk. The negative sign is what makes
the intuition work: the higher the payoff of the asset in bad times, the higher
$\operatorname{cov}(r_i,m)$ and $\beta_{i,m}$, and therefore the \emph{lower} its
expected return.}
\end{imfmlbox}

\begin{imkeybox}
Assets that pay off in bad times are valuable to hold, so prices for these assets
are high and expected returns are low. This is Lesson 6 stated as an equation.
\end{imkeybox}

% ------------------------------------------------------------------
\lo{IM-1 f}{Describe efficient market theory and explain how markets can be inefficient.}

\subsection{Efficient market theory}

\begin{itemize}
\item The ``classical'' notions of weak, semi-strong, and strong efficiency were
laid out by Fama (1970) and are obsolete.
\item In that year, the Nobel Prize committee also gave Robert Shiller the prize,
representing the opposite viewpoint (nearly efficient) of behavioral, or
non-rational, influences on financial markets.
\item Grossman and Stiglitz develop a model in which markets are
\term{near-efficient}. Active managers search for pockets of inefficiency, and in
doing so cause the market to be almost efficient.
\item In these pockets of inefficiency, active managers earn excess returns as a
reward for gathering and acting on costly information.
\end{itemize}

\begin{imdefbox}[title={The Grossman--Stiglitz paradox}]
Due to the existence of information costs, market efficiency and competitive
equilibrium are incompatible, and prices cannot be fully displayed. If prices
already revealed everything, nobody would pay to gather information --- and then
prices would reveal nothing.
\end{imdefbox}

\subsection{Why markets can be inefficient}

\begin{itemize}
\item In a \term{rational} explanation, high returns compensate for losses during
bad times. The key is defining those bad times and deciding whether these are
actually bad times for an individual investor.
\item In a \term{behavioral} explanation, high expected returns result from
agents' under- or overreaction to news or events. Behavioral biases can also
result from the inefficient updating of beliefs or ignoring some information.
\item For some risk premiums, the most compelling explanations are rational (as
with the volatility risk premium), for some behavioral (e.g.\ momentum), and for
some others a combination of rational and behavioral stories prevails (like
value/growth investing).
\end{itemize}

\begin{imtrapbox}[title={Consolidated trap summary --- IM-1}]
\begin{itemize}
\item \term{Sign, IM-1 c.} $E(R_m)-R_f = \bar{\gamma}\sigma_m^2$. Both terms push
the premium the same way. Higher risk aversion \emph{or} higher market variance
raises the premium.
\item \term{Sign, IM-1 e.} $b$ in $m = a + b\,r_m$ is negative. Low $m$ means bad
times in the CAPM.
\item \term{Polarity, IM-1 e.} Higher covariance with the SDF gives a
\emph{lower} risk premium, because $\lambda_m$ carries a minus sign. Assets that
pay off in bad times are expensive and earn little.
\item \term{Scope, IM-1 e.} An SDF requires only no arbitrage. The CAPM requires
far more. Do not swap these round.
\item \term{Sibling, IM-1 d.} Lesson 3 is the same under both models. Lessons 1,
4 and 5 are the ones that change wording.
\item \term{Role, IM-1 f.} Grossman--Stiglitz is about the \emph{cost of
information} making full efficiency self-defeating. It is not a behavioral
argument.
\item \term{Definition, IM-1 a.} Nutrients are good and factor risks are bad.
The analogy is deliberately asymmetric and the asymmetry is the testable point.
\end{itemize}
\end{imtrapbox}

% ==================================================================
